\chapter{Introduction}
    Financial derivatives are contracts whose payoff depends on the value of an underlying asset , such as an interest rate, an exchange rate, an equity price, or a commodity price [\cite{blyth2014introduction}]. There are many different types of financial derivatives contracts, such as options, forwards, futures and swaps. They are used extensively by financial institutions, corporations, and governments for hedging exposures, managing risk, and taking speculative positions \cite{hull2022options}. The accurate valuation of these instruments is therefore not merely a technical exercise but a practical necessity with direct implications for risk management, regulatory compliance, and financial stability \cite{bcbs2010basel3}. \\

    \noindent
    The standard approach to pricing financial derivatives rests on the principle of risk-neutral valuation, which states that the fair value of a derivative is the expected value of its future payoff, discounted at the risk-free rate, under a specially constructed probability measure \cite{black1973pricing, cox1976valuation}. This measure is derived from the requirement that no riskless profit can be extracted from the market \cite{harrison1979martingales, harrison1981martingales}. The elegance of this framework lies in its generality, once the risk-neutral measure is established, the price of any derivative reduces to computing an expectation, regardless of the specific form of the payoff. \\

    \noindent
    In practice, however, the classical risk-neutral framework rests on assumptions that are difficult to defend in real financial markets. It assumes that all market participants can borrow and lend freely at a common risk-free rate, that there are no costs associated with collateral or funding, and that the risk of counterparty default is negligible \cite{Gregory_2015}. The global financial crisis of 2007-2009 made clear that these were not merely simplifications but material omissions. The default of Lehman Brothers demonstrated that counterparty credit risk was real, that funding costs varied significantly across institutions, and that the assumption of a universal risk-free rate was an idealisation the market could no longer sustain \cite{Gregory_2015, Green_2016}. \\

    \noindent
    In response, the financial industry adopted a framework in which the economic value of a derivative is computed as the classical risk-free price adjusted by a set of additive corrections, collectively referred to as x-value adjustments (XVAs) \cite{Tsuchiya_2019}. Each adjustment captures a specific cost or risk that the classical framework had ignored. Credit Valuation Adjustment (CVA) accounts for the expected loss arising from counterparty default. Debit Valuation Adjustment (DVA) reflects the institution's own default risk as seen by its counterparty. Funding Valuation Adjustment (FVA) captures the cost of funding uncollateralised exposures \cite{Gregory_2015}. The framework has since been extended to include Margin Valuation Adjustment (MVA) for the cost of posting initial margin and Capital Valuation Adjustment (KVA) for the ongoing cost of holding regulatory capital \cite{Green_2016}. Together, these adjustments bring the theoretical price of a derivative closer to its true economic cost. \\ 

    \noindent
    Central to the computation of most XVA components is the concept of credit exposure \footnote{Credit exposure should not be confused with Value-at-Risk (VaR)}. Exposure measures the potential loss to a party in the event that its counterparty defaults at some future point in time \cite{Gregory_2015}. It is not a fixed quantity, as market conditions evolve, so does the value of the derivative and hence the potential loss. XVA calculations therefore require a complete picture of how a derivative's value may evolve over its entire lifetime under a wide range of market scenarios. This picture is captured by the exposure profile, specifically the expected positive exposure (EPE) and expected negative exposure (ENE), which average the positive and negative values of the derivative across simulated scenarios at each monitoring date \footnote{A monitoring date is a discrete future time point at which exposure is evaluated}. CVA and DVA are computed by integrating these profiles against default probabilities and survival functions, making accurate exposure profile estimation the central computational task in XVA valuation \cite{Gregory_2015}. \\

    \noindent
    Computing exposure profiles under classical methods requires simulating thousands of future market scenarios and re-pricing the derivative at each monitoring date along every simulated path \cite{deelstra2024accelerated}. This produces a nested simulation structure: an outer loop over market scenarios and an inner loop of re-pricing evaluations at each time step. The computational cost scales as the product of the number of outer paths and the cost of the inner re-pricing, making it expensive even for simple instruments \cite{Tsuchiya_2019}. For derivatives with complex payoff structures, path-dependent features, or long maturities, the cost becomes prohibitive. Large financial institutions manage portfolios of thousands of trades, each requiring its own exposure simulation, and must perform these calculations frequently for risk reporting and regulatory purposes \cite{KPMG2025DerivValuation}. Classical Monte Carlo methods are increasingly inadequate to meet this demand within the time constraints of real-world operations \cite{deelstra2024accelerated}. \\

    \noindent
    Deep learning surrogate models offer a practical approach to this problem. A surrogate is a neural network trained to approximate a function that is expensive to evaluate directly. In our context, the surrogate models will target the computation of the exposure profiles. Rather than simulating paths and re-pricing at each node, this study trains two deep learning surrogates to approximate the EPE and ENE profiles directly as functions of the initial yield curve, Hull-White model parameters, and monitoring time. Once trained, the surrogate can generate exposure profiles at a fraction of the cost of full Monte Carlo simulation. This technique is similar to \cite{frode2022neural}, in which \\
    
    \noindent
    Machine learning and deep learning have found broad application in quantitative finance as tools for overcoming the computational limitations of classical numerical methods. The core appeal is their capacity to approximate complex, high-dimensional functions at negligible inference cost once training is complete \cite{Goodfellow-et-al-2016}. In derivative pricing, this has proven particularly valuable for problems governed by high-dimensional partial differential equations, where classical finite difference and Monte Carlo methods suffer from a exponential growth in computational cost with the number of state variables \cite{han2018solving, beck2020overview}. The recognition that neural networks can serve as fast, accurate function approximators for such problems has driven a substantial body of research at the intersection of machine learning and computational finance \cite{weinan2020machine}. \\

    \noindent
    In the specific context of xVA computation, surrogate models have emerged as a practical solution to the nested simulation problem. Several studies have demonstrated that neural networks can be trained to approximate the mark-to-market value of an interest rate derivative across a range of market states, replacing the expensive inner re-pricing at each scenario node with a single forward pass \cite{frode2022neural, menegon2023real}. This approach has been extended to high-dimensional Bermudan options \cite{andersson2021deep}, to bilateral counterparty risk under stochastic intensity models \cite{gnoatto2022deep}, and to the computation of exposure and xVA profiles for portfolios of trades \cite{zhu2019applying}. The consistent finding across this literature is that well-designed surrogate architectures reproduce exposure profiles with high accuracy whilst reducing computation time substantially relative to full Monte Carlo simulation \cite{deelstra2024accelerated}. \\

    \noindent
    A limitation of standard surrogate training is that optimising on price or exposure labels alone does not guarantee that the network learns accurate sensitivities. In xVA applications, curve sensitivities are indispensable, they feed directly into hedging decisions and regulatory capital calculations, and a surrogate that prices correctly but produces incorrect gradients
    is of limited practical value \cite{deelstra2024accelerated}. Differential deep learning addresses this by incorporating analytic sensitivities as auxiliary training targets, optimizing the network jointly on values and their derivatives with respect to the input yield curve. This additional gradient information regularizes the learned mapping, improves generalization to sparsely sampled regions of the input space, and produces sensitivity estimates consistent with the underlying financial model \cite{lee2026deeponet}.\\

    \noindent
    This study applies deep learning to approximate the Expected Positive Exposure and Expected Negative Exposure profiles of an interest rate swap under the Hull-White one-factor model. Rather than training a surrogate to approximate the swap mark-to-market as an intermediate step, the networks are trained directly on EPE and ENE as functions of the initial yield curve, Hull-White parameters, and monitoring time. Since EPE and ENE are already expectations over the short rate distribution, the short rate does not enter the surrogate input, and CVA and DVA are recovered at inference through a single forward pass per monitoring date with no simulation. A mark-to-market surrogate is developed first as a validation step to confirm that the Hull-White pricing pipeline is correctly specified before the exposure targets are introduced. \\

    \noindent
    Standard surrogate training, which optimises on price or exposure labels alone, produces networks that price accurately in-distribution but reproduce sensitivities poorly. \cite{lee2026deeponet} demonstrate this directly for European bond options under the Hull-White and G2++ models, across all three surrogate paradigms evaluated, strong in-distribution pricing accuracy does not translate into reliable autodiff-based Vega estimates, with degradation most pronounced near at-the-money regions and under out-of-distribution volatility shifts. The reason is structural. Minimizing a price loss imposes no constraint on the gradient of the learned function, so the network is free to interpolate prices correctly while producing gradients that bear little resemblance to the true sensitivities of the underlying model. \\

    \noindent
    In the XVA setting, this matters considerably, because curve sensitivities feed directly into hedging decisions and regulatory capital calculations, and a surrogate that prices correctly but differentiates incorrectly is of limited practical value \cite{deelstra2024accelerated}. This gap between pricing accuracy and sensitivity fidelity motivates the use of differential training as the primary methodology in this study. \\


  % TODO: Results preview paragraph — write after training experiments are complete.
  % Cover: EPE/ENE accuracy vs MC benchmark, sensitivity fidelity (R^2 on dEPE/dy),
  % OOD robustness under volatility stress, and CVA/DVA propagation error.

    
\chapter{Background and Literature Review}
    \noindent
    The valuation of a financial derivative proceeds through a structured sequence of operations, beginning with the construction of market-consistent term structures from observable instrument prices and ending with the numerical evaluation of the contract's risk-neutral expectation under a calibrated stochastic model. The first stage, term structure construction, involves fitting yield curves, discount curves, or volatility surfaces to a cross-section of market quotes, a process that must respect no-arbitrage constraints whilst accommodating the irregularity and incompleteness of available market data \cite{nisha2024foundations}. \\

    \noindent
    Model calibration follows, in which the parameters of a chosen stochastic model are estimated by minimizing the discrepancy between model-implied and market-observed prices across a set of liquid benchmark instruments \cite{chang2024stochastic, choi2023calibrate}. Once the model is calibrated, each derivative in the portfolio must be priced individually, with the resulting values used to compute risk sensitivities, margin requirements, and valuation adjustments that feed into downstream risk management and reporting systems. The pipeline is therefore not a single calculation but a chain of dependent computations, each of which inherits the errors and computational costs of the stage preceding it. \\

    \noindent
    Each stage of this pipeline carries a high computational cost, and the burden compounds as model complexity and portfolio dimensionality increase. Model calibration requires repeated evaluation of the pricing function across many candidate parameter configurations, which is tractable for models with closed-form solutions but becomes expensive under stochastic volatility or interest rate models where each evaluation requires a numerical procedure \cite{horvath2020deeplearning}.  Pricing itself, for instruments without analytical solutions, typically demands either Monte Carlo simulation or a partial differential equation solver, both of which scale poorly with the
    dimension of the state space and the number of time steps required for accurate discretisation \cite{hutzenthaler2020overcoming, elbrachter2021dnn}.  \\

    \noindent
    The approximation of high-dimensional pricing functions via classical numerical methods is further constrained by the curse of dimensionality, whereby the computational cost of grid-based or quadrature-based methods grows exponentially with the number of state variables, rendering them impractical for multi-asset derivatives or multi-factor term structure models \cite{hutzenthaler2020overcoming}. The result is that even a single accurate valuation of a complex derivative can take minutes, which is acceptable in batch processing but incompatible with the response times required in real-time risk management.\\

    \noindent
    In a live trading risk management environment, this per-instrument computational burden is compounded by the scale and frequency at which valuations must be produced. A derivatives portfolio may contain thousands of instruments across multiple asset classes, and regulatory frameworks such as the Fundamental Review of the Trading Book require that risk sensitivities be computed and reported under many simultaneous market scenarios within an intraday timeframe \cite{despiegeleer2018machine}. Recalibration of the model is itself triggered by each significant market move, meaning that the full pipeline, from curve construction through to pricing and Greeks, must be re-executed many times within a single trading session \cite{anderson2023accelerated}. \\

    \noindent
    For instruments whose valuation depends on future exposure profiles, such as those subject to counterparty credit risk adjustments, the problem is further complicated by the need to simulate the portfolio's value across a grid of future dates and market states, a task that requires nested Monte Carlo simulation at a cost that scales with both portfolio size and the number of monitoring dates \cite{dhandapani2024optimizing}. This combination of per-instrument cost and portfolio-level scale creates a computational bottleneck that cannot be resolved by incremental improvements to existing numerical methods alone, and it is this bottleneck that has motivated the adoption of deep learning surrogates as a structural solution.\\

    \noindent
    Among the operations described above, the computation of xVAs stands out as the most computationally demanding, because it requires simulating the future value of every instrument in the portfolio across a grid of monitoring dates and risk scenarios, aggregating these values into exposure profiles, and then integrating those profiles against default intensity and loss assumptions to produce CVA and DVA estimates \cite{andersson2021deeplearning, goudenege2025xva}. The central computational objects in this pipeline are the Expected Positive Exposure and Expected Negative Exposure profiles, which must be evaluated at each monitoring date by taking expectations of the positive and negative parts of the portfolio value over the distribution of future market states, a task that requires nested Monte Carlo simulation when no closed-form expression is available. This study focuses specifically on this bottleneck, targeting EPE and ENE directly as surrogate outputs, and evaluating whether deep learning methods can reproduce these profiles accurately enough to replace full simulation in the CVA and DVA computation. \\

    \noindent
    The use of deep learning in derivatives pricing can be traced back to the seminar paper by \cite{hutchinson1994nonparametric}, who in 1994 proposed learning networks as a nonparametric alternative to the Black-Scholes formula for pricing and hedging equity options. By using data-driven methods like radial basis function network and multilayer perceptrons, the authors demonstrate that these networks can successfully approximate the Black-Scholes formula without requiring specific assumptions about underlying asset dynamics. The approach was validated by showing that the network-implied delta-hedging strategy achieved performance comparable to the Black-Scholes delta on out-of-sample S\&P 500 index options, despite the network not know the Black-Scholes formula. \\
    
    \noindent
    The motivation behind the seminal paper of Hutchinson \cite{hutchinson1994nonparametric} was the desire to avoid imposing parametric model assumptions in the Black-Scholes pricing framework. The dominant paradigm has since shifted, modern applications do not use neural networks because the reference model is distrusted, but because solving it accurately is
    too slow. The issue is not that the underlying mathematical formulations are incorrect. Rather, the governing stochastic differential equations and their associated pricing PDEs frequently do not admit closed-form solutions, and the numerical methods required to solve them, whether Monte Carlo simulation or finite difference schemes, often carry a computational cost that becomes prohibitive at the scale and frequency demanded by live risk management systems. \\

    \noindent
    For American and Bermudan options, where no closed-form solution is available, each pricing call requires Monte Carlo simulation or a backward induction procedure, and the need to repeat this across large portfolios under many scenarios creates a bottleneck that classical numerical methods cannot resolve within real-time constraints \cite{hutzenthaler2020overcoming, anderson2023accelerated}. The response in the recent literature has been to treat the neural network not as a model-free alternative but as a surrogate for a trusted reference model. The network is trained offline on synthetically generated labels and deployed online as an approximation to the model's own pricing function, preserving model-consistency whilst eliminating the per-query computational cost \cite{despiegeleer2018machine, becker2023learning}. \\

    \noindent
    This shift in motivation has a direct consequence for how surrogates are designed and evaluated. When the network stands in not for the model itself but for the expensive numerical procedure the model requires, whether Monte Carlo simulation, a finite difference PDE solver, or a nested simulation loop for exposure computation, the relevant criterion becomes accuracy relative to what that procedure would have produced. The model's stochastic equations and risk-neutral pricing framework remain trusted and intact; what the surrogate replaces is the act
    of solving them numerically at query time. This distinction matters because it means the surrogate must reproduce not only the model's prices but also its sensitivity structure, since downstream applications such as hedging and xVA computation depend on both, and a network that prices correctly but differentiates incorrectly provides only partial utility within the risk
    management pipeline. \\

    
    \noindent
    Under stochastic volatility models such as Heston and the rough Bergomi model, the calibration problem compounds this cost further. Fitting these models to an observed implied volatility surface requires minimizing the discrepancy between model-implied and market prices across a grid of strikes and maturities, which in turn demands repeated evaluation of the pricing function at each iteration of the optimization routine \cite{horvath2020deeplearning}. For rough volatility models in particular, where the variance process is driven by a fractional Brownian motion with no Markovian representation, even a single pricing call requires a full Monte Carlo simulation, making the calibration loop computationally intractable under real-time constraints \cite{becker2023learning}. The result is that model complexity and calibration frequency are in direct tension: the richer the model, the more accurately it captures observed market dynamics, but also the slower it is to solve, and the more frequently it needs to be recalibrated as the market moves. \\

    \noindent
    The curse of dimensionality manifests most severely in xVA computation, where the computational challenge is not a single expensive pricing call but the repeated evaluation of an entire portfolio across a grid of future market states. Simulating CVA and DVA requires an outer Monte Carlo loop that generates future market scenarios and, at each scenario and each monitoring date, a re-pricing of every instrument in the portfolio under the simulated state, a procedure known as nested simulation whose cost scales with the product of the number of outer paths, the number of monitoring dates, and the per-instrument pricing cost \cite{andersson2021deeplearning, goudenege2025xva}. \\

    \noindent
    For portfolios containing instruments whose pricing functions have no closed-form solutions, or are path dependent, such as Bermudan swaptions or callable bonds, the inner pricing step itself requires simulation or backward induction, making the full nested procedure computationally prohibitive for any portfolio of realistic size \cite{dhandapani2024optimizing}. This is the bottleneck that most directly motivates the surrogate modelling approach adopted in this study, since a surrogate that replaces the inner pricing step with a single forward pass through a trained network eliminates the nested structure entirely. \\

    \noindent
    For interest rate derivatives priced under multi-factor term structure models, the dimensionality of the state space creates an additional layer of computational difficulty. Grid-based PDE methods, which discretize the state space on a mesh, become impractical beyond two or three factors because the number of grid points grows exponentially with the number of dimensions, a phenomenon that Hutzenthaler et al.\ \cite{hutzenthaler2020overcoming} identify as a fundamental obstacle to classical numerical methods in high-dimensional settings. Practitioners are
    therefore forced to rely on Monte Carlo methods, which avoid the grid but converge slowly and require many paths to achieve acceptable accuracy for smooth payoffs, particularly when sensitivities must also be estimated \cite{choi2023calibrate}. \\

    \noindent
    The literature responding to these computational challenges has developed several distinct paradigms, each positioning the neural network differently in relation to the reference financial model. Physics-informed neural networks embed the pricing PDE residual directly into the training objective, enabling the network to learn a solution surface without requiring a labelled dataset of model prices \cite{bai2022pinn}. Deep backward stochastic differential equation methods reformulate the pricing problem via the Feynman-Kac correspondence and approximate the solution process using forward simulation combined with a neural network parameterisation \cite{yu2023backward}. Reinforcement learning and deep hedging approaches train agents to optimise hedging strategies directly from simulated or historical paths, bypassing the pricing equation altogether \cite{cao2020deephedging}. \\

    \noindent
    Generative models, including generative adversarial networks and variational autoencoders, have been applied to calibration and implied volatility surface generation, learning the geometry of the pricing problem from data rather than solving it analytically \cite{cuchiero2020gan}. This study operates within the supervised surrogate paradigm, in which a reference model provides exact labels through closed-form or numerical evaluation and the neural network is trained to approximate those labels across the relevant input space, with the trained surrogate deployed as a fast substitute for the reference computation at inference time. \\

    \noindent
    Deep neural networks substantially extended the scope of surrogate modelling in quantitative finance, enabling accurate approximation of pricing functions under stochastic volatility models, rough volatility models, and high-dimensional settings where classical numerical methods are computationally prohibitive \cite{horvath2020deeplearning, hutzenthaler2020overcoming}. Within the fixed-income and counterparty credit risk literature, feed-forward networks have been applied to CVA computation for portfolios of mixed derivatives, demonstrating that surrogate-based approaches can reproduce Monte Carlo exposure estimates with acceptable accuracy and considerable computational savings relative to full simulation \cite{andersson2021deeplearning}. \\

    \noindent
    For interest rate swap portfolios specifically, Vanhala \cite{vanhala2023dataxva} develops a data-driven approach to exposure profile computation using recurrent neural network architectures trained on simulated path data, showing that sequence models can learn exposure profiles without requiring repeated Monte Carlo evaluation at each monitoring date. Similarly, Dhandapani and Jain \cite{dhandapani2024optimizing} apply neural network surrogates to future exposure computation for Bermudan options, achieving significant reductions in inference cost relative to full backward induction. Across these contributions, the consistent finding is that deep surrogates generalise to fixed-income and exposure settings with strong in-distribution accuracy, but validation remains focused on price or exposure fidelity, with sensitivity accuracy receiving comparatively little attention. \\

    \noindent
    A structurally distinct class of surrogate models treats the pricing problem as learning an operator from a function space to a scalar output, replacing the finite-dimensional parameter
  vector with a yield curve or volatility surface as the network input. The deep operator network (DeepONet) is the canonical architecture in this setting, coupling a branch network that
  encodes the input function with a trunk network that encodes the query location, and was shown to approximate solution operators of parametric partial differential equations accurately
  across a broad range of input configurations \cite{wang2021learning}. For interest rate derivatives, this functional perspective is particularly natural: the pricing problem
  amounts to learning an operator from the space of initial term structures to bond or swap values across a range of contract and model parameters, and Cui et al.\
  \cite{cui2023learning} develop this idea for general derivative pricing maps under term structure models. \\

  \noindent
  Lee et al.\ \cite{lee2026deeponet} apply it directly to European bond call
  options under the Hull-White and G2++ models, using historical US Treasury yield curves as inputs and comparing DeepONet against PINN and DeepBSDE paradigms across pricing accuracy, Vega
  fidelity, and out-of-distribution robustness. Supervised operator learning achieves the strongest performance across all three criteria, but the authors identify a persistent gap: accurate prices do not automatically produce reliable autodiff-based sensitivities, particularly near the at-the-money region and under distribution shift. This finding establishes that the
  objective mismatch between price accuracy and sensitivity fidelity is not resolved by architectural sophistication alone, and motivates a training objective that targets sensitivities
  explicitly. \\

  \noindent
    A recurring finding across the surrogate literature is that minimising a price loss alone produces networks that price accurately in-distribution but reproduce sensitivities poorly, a gap with direct consequences for any application that depends on derivative outputs rather than prices alone. Huge and Savine \cite{huge2020differential} provide the foundational treatment of this problem through the differential machine learning framework, which trains neural networks jointly on price labels and their analytic derivatives with respect to inputs, exploiting the fact that analytic sensitivities are available at negligible additional cost whenever the reference model admits closed-form or semi-analytic expressions. The joint training objective regularises the learned function, concentrating it on the correct gradient structure rather than allowing the network to interpolate prices through arbitrary paths in function space. \\

    \noindent
    Subsequent work confirms the benefit of this approach across a range of settings: Mrad and Iassamen \cite{mrad2025differential} demonstrate that incorporating differential objectives accelerates convergence and improves sensitivity estimates for standard equity options, whilst the controlled comparisons of Lee et al.\ \cite{lee2026deeponet} for bond options under Hull-White show that derivative-aware training materially reduces the objective mismatch observed in purely price-supervised surrogates. Together, these contributions establish differential training as the appropriate methodology when sensitivity fidelity is a design requirement, and directly motivate its adoption in this study. \\

  \noindent
  Despite this body of work, a specific configuration remains largely unexamined: the direct use of differential deep learning surrogates to approximate expected exposure profiles for
  interest rate swaps, and the downstream propagation of surrogate-level accuracy into CVA and DVA estimates. Existing fixed-income surrogate studies either target mark-to-market values or,
  in the case of exposure-oriented work, use architectures trained on price or path labels without exploiting analytic sensitivities as training signals \cite{vanhala2023dataxva,
  dhandapani2024optimizing}. \\

  \noindent
  Work on CVA computation via neural networks has demonstrated that surrogate approaches can reproduce Monte Carlo CVA for portfolios of derivatives
  \cite{andersson2021deeplearning}, but the connection between the sensitivity fidelity of the surrogate and the accuracy of the resulting xVA estimates has not been examined systematically.
  This study addresses both gaps directly. Differential deep learning surrogates are trained with Expected Positive Exposure and Expected Negative Exposure as direct targets for a ten-year
  interest rate swap under the Hull-White one-factor model, and a downstream CVA propagation analysis translates surrogate-level accuracy metrics into xVA-level errors interpretable within a
  risk-management framework.

\chapter{Thesis Outline}
